% ID: SEL-D0074-OA-E03
% Tanda S074.
\textbf{Parte I.} Sean $A$ y $B$ los sucesos de resolver el primer y el segundo ejercicio, respectivamente.

\textbf{a)} El complementario de no resolver ninguno es resolver al menos uno, así que $P(A\cup B)=1-0{,}35=0{,}65$. Por inclusión y exclusión,
\[
 P(B)=0{,}65-0{,}30+0{,}10=\boxed{0{,}45}.
\]
\textbf{b)} Resolver solo el segundo tiene probabilidad $0{,}45-0{,}10=0{,}35$. Dentro del grupo que no resuelve el primero, de probabilidad $0{,}70$, obtenemos
\[
 P(B\mid A^c)=\frac{0{,}35}{0{,}70}=\boxed{0{,}50}.
\]

\textbf{Parte II. a)} La media muestral es $1800/9=200$ cm. La población es normal y $\sigma=4{,}5$ cm es conocida, así que usamos el intervalo normal para la media. Al $97\%$, cada cola tiene probabilidad $0{,}015$ y $z_{0{,}985}\approx2{,}1701$:
\[
 200\pm 2{,}1701\frac{4{,}5}{\sqrt9}
 \approx\boxed{[196{,}745;\ 203{,}255]}\ \text{cm}.
\]
\textbf{b)} Exigimos que $2{,}1701\cdot4{,}5/\sqrt n<1$. Al elevar al cuadrado resulta $n>(2{,}1701\cdot4{,}5)^2\approx95{,}37$. El menor entero que satisface la desigualdad estricta es $\boxed{96}$.
